\documentclass[../main.tex]{subfiles}
\begin{document}
\chapter{Probabilistic model}
\label{sec:general_series_system}
A probabilistic model is specified by equations involving random variables which make assumptions about how observable data about the system is generated. In what follows, we describe our model and observable data.

\section{Series system}
The system consists of $m$ components. We make the following assumption about 
the state of the components.
\begin{assumption}
\label{asm:binary_state}
The \jth component is initially in a non-failed state and permanently enters a failed state at some uncertain time $t_j > 0$ for $j=1\,,\ldots\,,m$.
\end{assumption}
Components have uncertain lifetimes as given by the following definition.
\begin{definition}
The lifetime of the \jth component is given by the continuous random variable 
$\cv_j \in \left(0\,,\infty\right)$ for $j=1\,,\ldots\,,m$.
\end{definition}
We make the following assumption about the joint distribution of the component lifetimes.
\begin{assumption}
\label{asm:indep}
The random variables $\cv_1\,, \ldots\,, \cv_m$ are mutually independent.
\end{assumption}

The system is a series system.
\begin{assumption}
    The system is in a non-failed state if all $m$ components are in a 
    non-failed state, otherwise the system is in a failed state.
\end{assumption}
A system in which at least $k$ of the components must be in a non-failed state 
for the system to be in a non-failed system is denoted a $k$-out-of-$m$ system, 
thus a series system is an $m$-out-of-$m$ system.

The uncertain lifetime of the system is given by the following definition.
\begin{theorem}
    \label{thm:series_system_lifetime}
    The lifetime of the system is given by the random variable
    \begin{equation}
    \label{eq:series_system_lifetime}
    \sv = \min{\!\left(\cv_1\,, \ldots\,, \cv_m\right)} \in \left(0, 
    \infty\right)\,,
    \end{equation}
    where $\cv_1\,, \ldots\,, \cv_m$ are the component lifetimes.
\end{theorem}
\begin{proof}
    A series system fails whenever any component fails, therefore its lifetime 
    is equal to the lifetime of the component with the minimum lifetime.
\end{proof}

\section{Probability distribution functions and parametric families}
As random variables, it is not certain which values $\sv,\cv_1,\ldots,\cv_m$ 
will realize. However, their respective probability distribution functions 
quantify the uncertainty. The cumulative distribution function is given by the 
following definition.
\begin{definition}
    \label{def:cdf_function}
    Let $\RV{X}$ be some random variable. The probability that $\RV{X} \leq x$ 
    is given by its cumulative distribution function, denoted by
    \begin{equation}
    \cdf_{\RV{X}}\!\left(x\right) = \Prob{\RV{X} \leq x}\,.
    \end{equation}
\end{definition}
The probability density function is given by the following definition.
\begin{definition}
    \label{def:pdf}
    Let $\RV{X}$ be some continuous random variable. The \emph{relative 
        likelihood} that $\RV{X} = x$ is given by its probability density 
    function, denoted by
    \begin{equation}
    \pdf_{\RV{X}}(x) = \derivative{t}\cdf_{\RV{X}}(t)\,.
    \end{equation}
\end{definition}
By the Second Fundamental Theorem of Calculus, the probability that a 
continuous random variable $\RV{X} \in \left(a,b\right]$, $a \leq b$, is given 
by
\begin{equation}
\label{eq:prob_F_R_p}
\Prob{a < \RV{X} \leq b} = \cdf_{\RV{X}}(b) - \cdf_{\RV{X}}(a) = \int_{a}^{b} 
\pdf_{\RV{X}}(s) \mathrm{d}s\,.
\end{equation}

By these definitions, the relative likelihood that the system will fail at time 
$t$ is $\spdf(t)$ and probability the \jth component will before time $t$ is 
$\cdf_{\cv_j}(t)$.
\begin{notation}
We simplify the notation and let $\cdf_j$ denote $\cdf_{\cv_j}$ and $\pdf_j$ denote $\pdf_{\cv_j}$.
Column vectors are denoted by boldface lower case letters, e.g., $\vec{x}$, and row vectors are denoted by the transpose of column vectors, e.g., $\vec{x}^\Transpose$.
Matrices are denoted by boldface upper case greek letters, e.g., $\matrx{A}$.
The \jth column of $\matrx{A}$ is denoted by $\matrx{A}_j$ (or $\left[\matrx{A}\right]_j$) and the $(j,k)$-th element of $\matrx{A}$ is denoted by $\matrx{A}_{j\,k}$ (or $\left[\matrx{A}\right]_{j\,k}$).
\end{notation}

We restrict our attention to parametrized families of probability distribution functions. The system lifetime has a probability density function that is a member of the following parametric family.
\begin{definition}
    \label{def:sysparametric}
    The system lifetime has a probability density function that is a member of 
    the parametric family of probability density functions denoted by
    \begin{equation}
    \label{eq:paramsys}
    \pfam = \left \{ \spdf\left(\; \cdot \given \sysparam\right) \colon 
    \sysparam \in \RealSet^{q \times m} \right \}\,,
    \end{equation}
    where $\sysparam_j \in \boldsymbol\Omega$ corresponds to the parameter 
    index of the \jth component lifetime.
\end{definition}
Particular values of $\sysparam$ index particular probability density functions 
within the $\pfam$-family. The true probability density function of the system 
lifetime is given by the following definition.
\begin{definition}
    \label{def:true_param}
    The distribution of the system lifetime has a true parameter index denoted 
    by $\tparam\sysparam$, and we say that
    \begin{equation}
        \sv \sim \spdf\!\left(\;\cdot \given \tparam\sysparam\right)\,.
    \end{equation}
\end{definition}
By these definitions, it follows that
\begin{equation}
    \cv_j \sim \pdf_j\!\left(\; \cdot \given \tparam\sysparam_j\right)\,.
\end{equation}

\section{Candidate selection models}
By \cref{thm:series_system_lifetime}, if a system failure occurs at time $\ti$, then one of the components failed at time $\ti$ (only a single component can cause a system failure since two or more continuous random variables cannot realize the same value).
\begin{definition}
\label{def:comp_cause}
	The discrete random variable indicating the component responsible for a system failure is denoted by $\RV{K}$ with a sample space $\Set{K} = \{1,\ldots,m\}$. That is, $\RV{K} = j$ indicates that $\cv_j < \cv_k$ for $k \in \SetDiff[\Set{K}][\{ j \}]$.
\end{definition}

We may not be certain which component caused a system failure, but may be able to narrow the component cause to some \emph{subset}.
The set of all component sets is given by the power set of $\Set{K}$, denoted by $\PS{\Set{K}}$.
The uncertain component sets are given by the following definition.
\begin{definition}
\label{def:random_candidate_set}
	Let $\rvcand$ denote the random set of components that are considered to be likely causes of a system failure. The sample space of $\rvcand$ is 
	$\PS{\Set{K}}$.
\end{definition}

The series system consists of $m$ components with random lifetimes $\cv_1,\ldots,\cv_m$.
These components are \emph{unobservable} causes of the system failure.
However, we suppose there is some process $A$ that has some knowledge about the components and which generates a set of candidate component causes, as depicted in \cref{fig:true_model}.
\begin{figure}
	\label{fig:true_model}
	\caption{Graphical model of random variables}
	\begin{center}
		% Graphic for TeX using PGF
% Title: /home/spinoza/Documents/repos/papers/series_system_estimation_from_masked_system_failure_time_samples/research/true_graph_model.dia
% Creator: Dia v0.97+git
% CreationDate: Sat Jul  6 06:56:22 2019
% For: spinoza
% \usepackage{tikz}
% The following commands are not supported in PSTricks at present
% We define them conditionally, so when they are implemented,
% this pgf file will use them.
\ifx\du\undefined
  \newlength{\du}
\fi
\setlength{\du}{15\unitlength}
\begin{tikzpicture}[even odd rule]
\pgftransformxscale{1.000000}
\pgftransformyscale{-1.000000}
\definecolor{dialinecolor}{rgb}{0.000000, 0.000000, 0.000000}
\pgfsetstrokecolor{dialinecolor}
\pgfsetstrokeopacity{1.000000}
\definecolor{diafillcolor}{rgb}{1.000000, 1.000000, 1.000000}
\pgfsetfillcolor{diafillcolor}
\pgfsetfillopacity{1.000000}
\pgfsetlinewidth{0.100000\du}
\pgfsetdash{{\pgflinewidth}{0.200000\du}}{0cm}
\pgfsetmiterjoin
\pgfsetbuttcap
\definecolor{dialinecolor}{rgb}{0.498039, 0.498039, 0.498039}
\pgfsetstrokecolor{dialinecolor}
\pgfsetstrokeopacity{1.000000}
\pgfpathmoveto{\pgfpoint{41.000000\du}{10.000000\du}}
\pgfpathcurveto{\pgfpoint{41.000000\du}{10.000000\du}}{\pgfpoint{41.000000\du}{8.000000\du}}{\pgfpoint{44.000000\du}{8.000000\du}}
\pgfpathcurveto{\pgfpoint{47.000000\du}{8.000000\du}}{\pgfpoint{48.000000\du}{8.000000\du}}{\pgfpoint{49.000000\du}{8.000000\du}}
\pgfpathcurveto{\pgfpoint{50.000000\du}{8.000000\du}}{\pgfpoint{51.000000\du}{8.000000\du}}{\pgfpoint{54.000000\du}{8.000000\du}}
\pgfpathcurveto{\pgfpoint{57.000000\du}{8.000000\du}}{\pgfpoint{57.000000\du}{10.000000\du}}{\pgfpoint{57.000000\du}{10.000000\du}}
\pgfpathcurveto{\pgfpoint{57.000000\du}{10.000000\du}}{\pgfpoint{57.000000\du}{12.000000\du}}{\pgfpoint{54.000000\du}{12.000000\du}}
\pgfpathcurveto{\pgfpoint{51.000000\du}{12.000000\du}}{\pgfpoint{50.000000\du}{12.000000\du}}{\pgfpoint{49.000000\du}{12.000000\du}}
\pgfpathcurveto{\pgfpoint{48.000000\du}{12.000000\du}}{\pgfpoint{47.000000\du}{12.000000\du}}{\pgfpoint{44.000000\du}{12.000000\du}}
\pgfpathcurveto{\pgfpoint{41.000000\du}{12.000000\du}}{\pgfpoint{41.000000\du}{10.000000\du}}{\pgfpoint{41.000000\du}{10.000000\du}}
\pgfpathclose
\pgfusepath{stroke}
% setfont left to latex
\definecolor{dialinecolor}{rgb}{0.000000, 0.000000, 0.000000}
\pgfsetstrokecolor{dialinecolor}
\pgfsetstrokeopacity{1.000000}
\definecolor{diafillcolor}{rgb}{0.000000, 0.000000, 0.000000}
\pgfsetfillcolor{diafillcolor}
\pgfsetfillopacity{1.000000}
\node[anchor=base west,inner sep=0pt,outer sep=0pt,color=dialinecolor] at (43.000000\du,10.000000\du){};
% setfont left to latex
\definecolor{dialinecolor}{rgb}{0.000000, 0.000000, 0.000000}
\pgfsetstrokecolor{dialinecolor}
\pgfsetstrokeopacity{1.000000}
\definecolor{diafillcolor}{rgb}{0.000000, 0.000000, 0.000000}
\pgfsetfillcolor{diafillcolor}
\pgfsetfillopacity{1.000000}
\node[anchor=base west,inner sep=0pt,outer sep=0pt,color=dialinecolor] at (43.000000\du,10.000000\du){};
% setfont left to latex
\definecolor{dialinecolor}{rgb}{0.000000, 0.000000, 0.000000}
\pgfsetstrokecolor{dialinecolor}
\pgfsetstrokeopacity{1.000000}
\definecolor{diafillcolor}{rgb}{0.000000, 0.000000, 0.000000}
\pgfsetfillcolor{diafillcolor}
\pgfsetfillopacity{1.000000}
\node[anchor=base west,inner sep=0pt,outer sep=0pt,color=dialinecolor] at (45.013631\du,4.000000\du){};
\pgfsetlinewidth{0.100000\du}
\pgfsetdash{}{0pt}
\pgfsetbuttcap
\pgfsetmiterjoin
\pgfsetlinewidth{0.100000\du}
\pgfsetbuttcap
\pgfsetmiterjoin
\pgfsetdash{}{0pt}
\definecolor{dialinecolor}{rgb}{0.000000, 0.639216, 0.000000}
\pgfsetstrokecolor{dialinecolor}
\pgfsetstrokeopacity{1.000000}
\pgfpathmoveto{\pgfpoint{45.013631\du}{3.000000\du}}
\pgfpathcurveto{\pgfpoint{45.563631\du}{3.000000\du}}{\pgfpoint{46.013631\du}{3.450000\du}}{\pgfpoint{46.013631\du}{4.000000\du}}
\pgfpathcurveto{\pgfpoint{46.013631\du}{4.550000\du}}{\pgfpoint{45.563631\du}{5.000000\du}}{\pgfpoint{45.013631\du}{5.000000\du}}
\pgfpathcurveto{\pgfpoint{44.463631\du}{5.000000\du}}{\pgfpoint{44.013631\du}{4.550000\du}}{\pgfpoint{44.013631\du}{4.000000\du}}
\pgfpathcurveto{\pgfpoint{44.013631\du}{3.450000\du}}{\pgfpoint{44.463631\du}{3.000000\du}}{\pgfpoint{45.013631\du}{3.000000\du}}
\pgfusepath{stroke}
% setfont left to latex
\definecolor{dialinecolor}{rgb}{0.000000, 0.000000, 0.000000}
\pgfsetstrokecolor{dialinecolor}
\pgfsetstrokeopacity{1.000000}
\definecolor{diafillcolor}{rgb}{0.000000, 0.000000, 0.000000}
\pgfsetfillcolor{diafillcolor}
\pgfsetfillopacity{1.000000}
\node[anchor=base west,inner sep=0pt,outer sep=0pt,color=dialinecolor] at (44.739031\du,4.371370\du){S};
% setfont left to latex
\definecolor{dialinecolor}{rgb}{0.000000, 0.000000, 0.000000}
\pgfsetstrokecolor{dialinecolor}
\pgfsetstrokeopacity{1.000000}
\definecolor{diafillcolor}{rgb}{0.000000, 0.000000, 0.000000}
\pgfsetfillcolor{diafillcolor}
\pgfsetfillopacity{1.000000}
\node[anchor=base west,inner sep=0pt,outer sep=0pt,color=dialinecolor] at (49.000000\du,4.000000\du){};
% setfont left to latex
\definecolor{dialinecolor}{rgb}{0.000000, 0.000000, 0.000000}
\pgfsetstrokecolor{dialinecolor}
\pgfsetstrokeopacity{1.000000}
\definecolor{diafillcolor}{rgb}{0.000000, 0.000000, 0.000000}
\pgfsetfillcolor{diafillcolor}
\pgfsetfillopacity{1.000000}
\node[anchor=base west,inner sep=0pt,outer sep=0pt,color=dialinecolor] at (49.000000\du,4.000000\du){};
% setfont left to latex
\definecolor{dialinecolor}{rgb}{0.000000, 0.000000, 0.000000}
\pgfsetstrokecolor{dialinecolor}
\pgfsetstrokeopacity{1.000000}
\definecolor{diafillcolor}{rgb}{0.000000, 0.000000, 0.000000}
\pgfsetfillcolor{diafillcolor}
\pgfsetfillopacity{1.000000}
\node[anchor=base west,inner sep=0pt,outer sep=0pt,color=dialinecolor] at (49.000000\du,4.000000\du){};
\pgfsetlinewidth{0.100000\du}
\pgfsetdash{}{0pt}
\pgfsetbuttcap
\pgfsetmiterjoin
\pgfsetlinewidth{0.100000\du}
\pgfsetbuttcap
\pgfsetmiterjoin
\pgfsetdash{}{0pt}
\definecolor{dialinecolor}{rgb}{0.000000, 0.639216, 0.000000}
\pgfsetstrokecolor{dialinecolor}
\pgfsetstrokeopacity{1.000000}
\pgfpathmoveto{\pgfpoint{49.013631\du}{3.000000\du}}
\pgfpathcurveto{\pgfpoint{49.563631\du}{3.000000\du}}{\pgfpoint{50.013631\du}{3.450000\du}}{\pgfpoint{50.013631\du}{4.000000\du}}
\pgfpathcurveto{\pgfpoint{50.013631\du}{4.550000\du}}{\pgfpoint{49.563631\du}{5.000000\du}}{\pgfpoint{49.013631\du}{5.000000\du}}
\pgfpathcurveto{\pgfpoint{48.463631\du}{5.000000\du}}{\pgfpoint{48.013631\du}{4.550000\du}}{\pgfpoint{48.013631\du}{4.000000\du}}
\pgfpathcurveto{\pgfpoint{48.013631\du}{3.450000\du}}{\pgfpoint{48.463631\du}{3.000000\du}}{\pgfpoint{49.013631\du}{3.000000\du}}
\pgfusepath{stroke}
% setfont left to latex
\definecolor{dialinecolor}{rgb}{0.000000, 0.000000, 0.000000}
\pgfsetstrokecolor{dialinecolor}
\pgfsetstrokeopacity{1.000000}
\definecolor{diafillcolor}{rgb}{0.000000, 0.000000, 0.000000}
\pgfsetfillcolor{diafillcolor}
\pgfsetfillopacity{1.000000}
\node[anchor=base west,inner sep=0pt,outer sep=0pt,color=dialinecolor] at (48.602431\du,4.350340\du){C};
% setfont left to latex
\definecolor{dialinecolor}{rgb}{0.000000, 0.000000, 0.000000}
\pgfsetstrokecolor{dialinecolor}
\pgfsetstrokeopacity{1.000000}
\definecolor{diafillcolor}{rgb}{0.000000, 0.000000, 0.000000}
\pgfsetfillcolor{diafillcolor}
\pgfsetfillopacity{1.000000}
\node[anchor=base west,inner sep=0pt,outer sep=0pt,color=dialinecolor] at (50.000000\du,4.000000\du){};
% setfont left to latex
\definecolor{dialinecolor}{rgb}{0.000000, 0.000000, 0.000000}
\pgfsetstrokecolor{dialinecolor}
\pgfsetstrokeopacity{1.000000}
\definecolor{diafillcolor}{rgb}{0.000000, 0.000000, 0.000000}
\pgfsetfillcolor{diafillcolor}
\pgfsetfillopacity{1.000000}
\node[anchor=base west,inner sep=0pt,outer sep=0pt,color=dialinecolor] at (53.000000\du,4.000000\du){};
% setfont left to latex
\definecolor{dialinecolor}{rgb}{0.000000, 0.000000, 0.000000}
\pgfsetstrokecolor{dialinecolor}
\pgfsetstrokeopacity{1.000000}
\definecolor{diafillcolor}{rgb}{0.000000, 0.000000, 0.000000}
\pgfsetfillcolor{diafillcolor}
\pgfsetfillopacity{1.000000}
\node[anchor=base west,inner sep=0pt,outer sep=0pt,color=dialinecolor] at (54.400000\du,10.400000\du){Tm};
% setfont left to latex
\definecolor{dialinecolor}{rgb}{0.000000, 0.000000, 0.000000}
\pgfsetstrokecolor{dialinecolor}
\pgfsetstrokeopacity{1.000000}
\definecolor{diafillcolor}{rgb}{0.000000, 0.000000, 0.000000}
\pgfsetfillcolor{diafillcolor}
\pgfsetfillopacity{1.000000}
\node[anchor=base west,inner sep=0pt,outer sep=0pt,color=dialinecolor] at (55.000000\du,10.000000\du){};
\pgfsetlinewidth{0.100000\du}
\pgfsetdash{{\pgflinewidth}{0.200000\du}}{0cm}
\pgfsetmiterjoin
\pgfsetbuttcap
{
\definecolor{diafillcolor}{rgb}{0.000000, 0.000000, 0.000000}
\pgfsetfillcolor{diafillcolor}
\pgfsetfillopacity{1.000000}
% was here!!!
\pgfsetarrowsend{latex}
\definecolor{dialinecolor}{rgb}{0.000000, 0.000000, 0.000000}
\pgfsetstrokecolor{dialinecolor}
\pgfsetstrokeopacity{1.000000}
\pgfpathmoveto{\pgfpoint{49.000000\du}{8.000000\du}}
\pgfpathcurveto{\pgfpoint{49.000000\du}{6.000000\du}}{\pgfpoint{45.013631\du}{7.150000\du}}{\pgfpoint{45.013631\du}{5.150000\du}}
\pgfusepath{stroke}
}
\pgfsetlinewidth{0.100000\du}
\pgfsetdash{}{0pt}
\pgfsetbuttcap
\pgfsetmiterjoin
\pgfsetlinewidth{0.100000\du}
\pgfsetbuttcap
\pgfsetmiterjoin
\pgfsetdash{}{0pt}
\definecolor{dialinecolor}{rgb}{0.803922, 0.000000, 0.000000}
\pgfsetstrokecolor{dialinecolor}
\pgfsetstrokeopacity{1.000000}
\pgfpathellipse{\pgfpoint{46.000000\du}{10.000000\du}}{\pgfpoint{1.000000\du}{0\du}}{\pgfpoint{0\du}{1.000000\du}}
\pgfusepath{stroke}
% setfont left to latex
\definecolor{dialinecolor}{rgb}{0.000000, 0.000000, 0.000000}
\pgfsetstrokecolor{dialinecolor}
\pgfsetstrokeopacity{1.000000}
\definecolor{diafillcolor}{rgb}{0.000000, 0.000000, 0.000000}
\pgfsetfillcolor{diafillcolor}
\pgfsetfillopacity{1.000000}
\node[anchor=base west,inner sep=0pt,outer sep=0pt,color=dialinecolor] at (45.400000\du,10.400000\du){T2};
% setfont left to latex
\definecolor{dialinecolor}{rgb}{0.000000, 0.000000, 0.000000}
\pgfsetstrokecolor{dialinecolor}
\pgfsetstrokeopacity{1.000000}
\definecolor{diafillcolor}{rgb}{0.000000, 0.000000, 0.000000}
\pgfsetfillcolor{diafillcolor}
\pgfsetfillopacity{1.000000}
\node[anchor=base west,inner sep=0pt,outer sep=0pt,color=dialinecolor] at (46.000000\du,10.000000\du){};
\pgfsetlinewidth{0.100000\du}
\pgfsetdash{}{0pt}
\pgfsetbuttcap
\pgfsetmiterjoin
\pgfsetlinewidth{0.100000\du}
\pgfsetbuttcap
\pgfsetmiterjoin
\pgfsetdash{}{0pt}
\definecolor{dialinecolor}{rgb}{0.803922, 0.000000, 0.000000}
\pgfsetstrokecolor{dialinecolor}
\pgfsetstrokeopacity{1.000000}
\pgfpathellipse{\pgfpoint{43.000000\du}{10.000000\du}}{\pgfpoint{1.000000\du}{0\du}}{\pgfpoint{0\du}{1.000000\du}}
\pgfusepath{stroke}
% setfont left to latex
\definecolor{dialinecolor}{rgb}{0.000000, 0.000000, 0.000000}
\pgfsetstrokecolor{dialinecolor}
\pgfsetstrokeopacity{1.000000}
\definecolor{diafillcolor}{rgb}{0.000000, 0.000000, 0.000000}
\pgfsetfillcolor{diafillcolor}
\pgfsetfillopacity{1.000000}
\node[anchor=base west,inner sep=0pt,outer sep=0pt,color=dialinecolor] at (42.400000\du,10.400000\du){T1};
\pgfsetlinewidth{0.100000\du}
\pgfsetdash{}{0pt}
\pgfsetbuttcap
\pgfsetmiterjoin
\pgfsetlinewidth{0.100000\du}
\pgfsetbuttcap
\pgfsetmiterjoin
\pgfsetdash{}{0pt}
\definecolor{dialinecolor}{rgb}{0.803922, 0.000000, 0.000000}
\pgfsetstrokecolor{dialinecolor}
\pgfsetstrokeopacity{1.000000}
\pgfpathellipse{\pgfpoint{55.000000\du}{10.000000\du}}{\pgfpoint{1.000000\du}{0\du}}{\pgfpoint{0\du}{1.000000\du}}
\pgfusepath{stroke}
% setfont left to latex
\definecolor{dialinecolor}{rgb}{0.000000, 0.000000, 0.000000}
\pgfsetstrokecolor{dialinecolor}
\pgfsetstrokeopacity{1.000000}
\definecolor{diafillcolor}{rgb}{0.000000, 0.000000, 0.000000}
\pgfsetfillcolor{diafillcolor}
\pgfsetfillopacity{1.000000}
\node[anchor=base west,inner sep=0pt,outer sep=0pt,color=dialinecolor] at (55.000000\du,10.000000\du){};
\pgfsetlinewidth{0.100000\du}
\pgfsetdash{{\pgflinewidth}{0.200000\du}}{0cm}
\pgfsetbuttcap
{
\definecolor{diafillcolor}{rgb}{0.000000, 0.000000, 0.000000}
\pgfsetfillcolor{diafillcolor}
\pgfsetfillopacity{1.000000}
% was here!!!
\definecolor{dialinecolor}{rgb}{0.000000, 0.000000, 0.000000}
\pgfsetstrokecolor{dialinecolor}
\pgfsetstrokeopacity{1.000000}
\draw (47.500000\du,10.000000\du)--(53.451200\du,10.000000\du);
}
% setfont left to latex
\definecolor{dialinecolor}{rgb}{0.000000, 0.000000, 0.000000}
\pgfsetstrokecolor{dialinecolor}
\pgfsetstrokeopacity{1.000000}
\definecolor{diafillcolor}{rgb}{0.000000, 0.000000, 0.000000}
\pgfsetfillcolor{diafillcolor}
\pgfsetfillopacity{1.000000}
\node[anchor=base west,inner sep=0pt,outer sep=0pt,color=dialinecolor] at (53.000000\du,4.000000\du){};
% setfont left to latex
\definecolor{dialinecolor}{rgb}{0.000000, 0.000000, 0.000000}
\pgfsetstrokecolor{dialinecolor}
\pgfsetstrokeopacity{1.000000}
\definecolor{diafillcolor}{rgb}{0.000000, 0.000000, 0.000000}
\pgfsetfillcolor{diafillcolor}
\pgfsetfillopacity{1.000000}
\node[anchor=base west,inner sep=0pt,outer sep=0pt,color=dialinecolor] at (52.842360\du,3.886212\du){};
\pgfsetlinewidth{0.100000\du}
\pgfsetdash{}{0pt}
\pgfsetbuttcap
\pgfsetmiterjoin
\pgfsetlinewidth{0.100000\du}
\pgfsetbuttcap
\pgfsetmiterjoin
\pgfsetdash{}{0pt}
\definecolor{dialinecolor}{rgb}{0.803922, 0.000000, 0.000000}
\pgfsetstrokecolor{dialinecolor}
\pgfsetstrokeopacity{1.000000}
\pgfpathellipse{\pgfpoint{53.007601\du}{4.000000\du}}{\pgfpoint{1.000000\du}{0\du}}{\pgfpoint{0\du}{1.000000\du}}
\pgfusepath{stroke}
% setfont left to latex
\definecolor{dialinecolor}{rgb}{0.000000, 0.000000, 0.000000}
\pgfsetstrokecolor{dialinecolor}
\pgfsetstrokeopacity{1.000000}
\definecolor{diafillcolor}{rgb}{0.000000, 0.000000, 0.000000}
\pgfsetfillcolor{diafillcolor}
\pgfsetfillopacity{1.000000}
\node[anchor=base west,inner sep=0pt,outer sep=0pt,color=dialinecolor] at (52.596401\du,4.350340\du){K};
% setfont left to latex
\definecolor{dialinecolor}{rgb}{0.000000, 0.000000, 0.000000}
\pgfsetstrokecolor{dialinecolor}
\pgfsetstrokeopacity{1.000000}
\definecolor{diafillcolor}{rgb}{0.000000, 0.000000, 0.000000}
\pgfsetfillcolor{diafillcolor}
\pgfsetfillopacity{1.000000}
\node[anchor=base west,inner sep=0pt,outer sep=0pt,color=dialinecolor] at (52.787834\du,4.036156\du){};
% setfont left to latex
\definecolor{dialinecolor}{rgb}{0.000000, 0.000000, 0.000000}
\pgfsetstrokecolor{dialinecolor}
\pgfsetstrokeopacity{1.000000}
\definecolor{diafillcolor}{rgb}{0.000000, 0.000000, 0.000000}
\pgfsetfillcolor{diafillcolor}
\pgfsetfillopacity{1.000000}
\node[anchor=base west,inner sep=0pt,outer sep=0pt,color=dialinecolor] at (52.787834\du,3.899843\du){};
\pgfsetlinewidth{0.100000\du}
\pgfsetdash{{\pgflinewidth}{0.200000\du}}{0cm}
\pgfsetmiterjoin
\pgfsetbuttcap
{
\definecolor{diafillcolor}{rgb}{0.000000, 0.000000, 0.000000}
\pgfsetfillcolor{diafillcolor}
\pgfsetfillopacity{1.000000}
% was here!!!
\pgfsetarrowsend{latex}
\definecolor{dialinecolor}{rgb}{0.000000, 0.000000, 0.000000}
\pgfsetstrokecolor{dialinecolor}
\pgfsetstrokeopacity{1.000000}
\pgfpathmoveto{\pgfpoint{49.000000\du}{8.000000\du}}
\pgfpathcurveto{\pgfpoint{49.000000\du}{6.000000\du}}{\pgfpoint{49.000000\du}{7.150000\du}}{\pgfpoint{49.000000\du}{5.150000\du}}
\pgfusepath{stroke}
}
\pgfsetlinewidth{0.100000\du}
\pgfsetdash{{\pgflinewidth}{0.200000\du}}{0cm}
\pgfsetmiterjoin
\pgfsetbuttcap
{
\definecolor{diafillcolor}{rgb}{0.000000, 0.000000, 0.000000}
\pgfsetfillcolor{diafillcolor}
\pgfsetfillopacity{1.000000}
% was here!!!
\pgfsetarrowsend{latex}
\definecolor{dialinecolor}{rgb}{0.000000, 0.000000, 0.000000}
\pgfsetstrokecolor{dialinecolor}
\pgfsetstrokeopacity{1.000000}
\pgfpathmoveto{\pgfpoint{49.000000\du}{8.000000\du}}
\pgfpathcurveto{\pgfpoint{49.000000\du}{6.000000\du}}{\pgfpoint{53.007601\du}{7.150000\du}}{\pgfpoint{53.007601\du}{5.150000\du}}
\pgfusepath{stroke}
}
% setfont left to latex
\definecolor{dialinecolor}{rgb}{0.000000, 0.000000, 0.000000}
\pgfsetstrokecolor{dialinecolor}
\pgfsetstrokeopacity{1.000000}
\definecolor{diafillcolor}{rgb}{0.000000, 0.000000, 0.000000}
\pgfsetfillcolor{diafillcolor}
\pgfsetfillopacity{1.000000}
\node[anchor=base west,inner sep=0pt,outer sep=0pt,color=dialinecolor] at (47.049059\du,5.576493\du){A};
% setfont left to latex
\definecolor{dialinecolor}{rgb}{0.000000, 0.000000, 0.000000}
\pgfsetstrokecolor{dialinecolor}
\pgfsetstrokeopacity{1.000000}
\definecolor{diafillcolor}{rgb}{0.000000, 0.000000, 0.000000}
\pgfsetfillcolor{diafillcolor}
\pgfsetfillopacity{1.000000}
\node[anchor=base west,inner sep=0pt,outer sep=0pt,color=dialinecolor] at (50.375095\du,5.440180\du){};
\pgfsetlinewidth{0.050000\du}
\pgfsetdash{{\pgflinewidth}{0.200000\du}}{0cm}
\pgfsetmiterjoin
\pgfsetbuttcap
{
\definecolor{diafillcolor}{rgb}{0.313726, 0.313726, 0.313726}
\pgfsetfillcolor{diafillcolor}
\pgfsetfillopacity{1.000000}
% was here!!!
\pgfsetarrowsend{stealth}
\definecolor{dialinecolor}{rgb}{0.313726, 0.313726, 0.313726}
\pgfsetstrokecolor{dialinecolor}
\pgfsetstrokeopacity{1.000000}
\pgfpathmoveto{\pgfpoint{47.663599\du}{5.395892\du}}
\pgfpathcurveto{\pgfpoint{48.658683\du}{5.395892\du}}{\pgfpoint{47.895331\du}{6.200138\du}}{\pgfpoint{48.985834\du}{6.241032\du}}
\pgfusepath{stroke}
}
\end{tikzpicture}

	\end{center}
\end{figure}

We consider two ways failed component candidate sets may be constructed.
Suppose candidate sets are of cardinality $w$.
The first way candidate sets are generated is \emph{optimal} since we simply draw from the component lifetimes and include the $w$ components with the smallest lifetimes in the set. The probability that a component in a set $\Set{A} \in \PS{\Set{K}}$ is the cause of a system failure follows from the definition of $\RV{K}$, i.e., $\Prob{\RV{K} \in \Set{A}}$.

The second way is perhaps more realistic in that it is \emph{certain} that one of the components in the candidate set is in a failed state (and is thus the cause of the system failure), e.g., a diagnostician was able to isolate the problem to some subset of the system.
However, the $w-1$ non-failed componets in the candidate set are \emph{randomly} selected from the remaining $m-1$ components, which is probably overly pessimistic.
These two approaches generate a range, from useful but pessimistic to unreasonably \emph{optimal}. For now, we provide notation for the \emph{random component sets} and treat the two approaches separately in \cref{sec:opt,sec:pess}.


The \emph{cardinality} of a component set could theoretically convey 
information about the system. As a simplification, we make the following 
assumption.
\begin{assumption}
\label{asm:indep_W_S}
The cardinality of the random set $\rvcand$, denoted by the discrete random variable
\begin{equation}
\RV{W} = \Card{\rvcand}\,,
\end{equation}
is independent of the random system lifetime $\sv$ and the random component 
cause $\rvcand$ and therefore does not carry information about the true 
parameter index $\tparam\sysparam$.
\end{assumption}

The indicator function and the $k$-subset functions are needed for subsequent material and are given by the following definitions.
\begin{definition}[Indicator function]
	Given a subset $\Set{A}$ of a set $\Set{B}$, the indicator function $\SetIndicator{\Set{A}} \colon \Set{B} \mapsto \{0,1\}$
	is equal to $1$ if $x \in \Set{A}$ and $0$ otherwise.
\end{definition}

\begin{definition}
	The subsets of $\Set{A}$ that are of cardinality $k$ is defined as
	\begin{equation}
	\left[\Set{A}\right]^k \coloneqq \SetBuilder{\Set{X} \in \Set{A}}{\Card{\Set{A}} = k}\,,
	\end{equation}
	which has a cardinality
	\begin{equation}
	\binom{\Card{\Set{A}}}{k}\,.
	\end{equation}
\end{definition}

We are primarily interested in the distribution of candidate sets on the condition that they have a specified cardinality.
\begin{definition}
	The random variable $\rvcand$ conditioned on $\RV{W} = w$ has a sample space given by
	\begin{equation}
	\candsw{w} = \left[\PS{\RV{K}}\right]^w\,.
	\end{equation}
\end{definition}
The reason why a candidate set of cardinality $w$ is generated is outside the scope of the statistical model and where necessary, the relative frequency in a sample is used.

\section{Masked system failure times}
\begin{definition}
A \emph{masked system failure time} consists of a system failure time $\ti$ and a corresponding set of components $\cand$.
\end{definition}
\begin{assumption}
\label{asm:cand_prob}
A random \emph{masked system failure time} is a jointly distributed random 
system lifetime $\sv$ and random component set $\rvcand$.
\end{assumption}

\begin{assumption}
\label{def:msfs}
A random sample of $n$ \emph{masked system failure times}, denoted by
\begin{equation}
\mfs = \left( \Pair{\sv_1}{\rvcand_1},\ldots,\Pair{\sv_n}{\rvcand_n}\right)\,,
\end{equation}
consists of $n$ independent and identically jointly distributed pairs of system failure times and candidate sets.
\end{assumption}

\begin{assumption}
The only information we have about the system and its true parameter index 
$\tparam\sysparam$ is given by observing a random sample of $n$ masked system 
failure times,
\begin{equation}
    \left(
    	\Pair{\sv_1=\ti_1}{\rvcand_1=\Set{C}[1]}, \ldots, 
    	\Pair{\sv_n=\ti_n}{\rvcand_n=\Set{C}[n]}
    \right)\,.
\end{equation}
\end{assumption}

Given the assumed model, the object of statistical interest is the (unknown) true parameter index $\tparam\sysparam$.
Provided an estimate of $\tparam\sysparam$, any characteristic that is a function of $\tparam\sysparam$ can be estimated as a result, thus the primary objective is to use the \emph{information} in a sample of $n$ masked system failure times to estimate $\tparam\sysparam$ with some quantifiable uncertainty.

The likelihood of observing a particular realization, $\mfsw = \mfso$, is given by the following definition.
\begin{definition}
	\label{def:msft_joint_pdf}
	A random sample of $n$ masked system failure times conditioned on candidate sets of cardinality $w$ has a joint probability density given by
	\begin{equation}
	\label{eq:msft_joint_pdf}
	\pdf_{\mfsw}\left(\mfso \given w, \tparam{\sysparam}\right) =
	\prod_{i=1}^{n} \spdf(t_i \given \tparam\sysparam) \cdot \pmf_{\rvcand \given \sv, \RV{W}}\left(\cand_i \given t_i, w, \tparam\sysparam\right)\,,
	\end{equation}
	where $\spdf(\cdot \given \tparam\sysparam)$ is the marginal density given by \cref{thm:series_pdf_function} and $\pmf_{\rvcand \given \sv, \RV{W}}(\cdot \given t, w,, \tparam\sysparam)$ is the conditional probability mass given by \cref{def:general_cond_cand_probability}.
\end{definition}

The generative model of $\mfsw$, which directly follows from \cref{eq:msft_joint_pdf}, is given by \cref{alg:Fn_generator}.
\begin{algorithm}[h]
	\DontPrintSemicolon
	\KwResult{a realization of a masked system failure time with $w$ candidates.}
	\KwIn{\\
		$\qquad\tparam\sysparam$, the true parameter value.\\
		$\qquad w$, the cardinality of the candidate sets.}
	\KwOut{\\
		$\qquad \mfo$, a realization of $\rvcand, \sv$ conditioned on $\RV{W} = w$.}
	\BlankLine
	\SetKwProg{func}{Model}{}{}
	\func{\genSampleFn{$w, \tparam\sysparam$}}
	{
		draw a system lifetime $t \sim \spdf(\;\cdot \given \sysparam)$\;
		draw a candidate set $\cand \sim \pmf_{\rvcand \given \sv, \RV{W}}(\; \cdot \given t, w, \sysparam)$\;
		$\mfo \gets \Pair{t}{\cand}$\;
		\KwRet{$\mfo$}
	}
	\caption{Generative model of a masked system failure time with $w$ candidates.}
	\label{alg:Fn_generator}
\end{algorithm}
\end{document}